\section{Toolkit and Diagnostics}

\tool treats a tokenizer export as an artifact, not as an opaque component of a
trained recommender. The minimum input is a table
\texttt{sid\_assignments(item\_id, sid\_0, ..., sid\_L)}. Optional tables
provide item metadata, interaction histories, and generator outputs. This
contract lets adapters normalize heterogeneous repositories while keeping the
audit independent of a particular training loop.

\paragraph{D1: utilization.}
\tool reports per-level usage, prefix counts, entropy-style imbalance
summaries, and dead or underused code indicators. These expose codebook
collapse and uneven capacity even when full ranking metrics are unavailable.

\paragraph{D2: collision profile.}
\tool measures duplicate full \sid{}s and prefix collisions at each depth. In
the current resource version, D2 is a collision profile rather than a causal
harm estimate. Interaction-qualified collision harm, as motivated by
qualification-aware collision work~\cite{hu2026quasid}, is reserved for the
next artifact version.

\paragraph{D3: semantic-collaborative alignment.}
Semantic grouping and collaborative usefulness need not agree. \tool therefore
computes a co-occurrence-based collaborative neighborhood proxy and asks how
often prefix neighborhoods recover those behaviorally related items. Metadata
purity is retained only as an auxiliary diagnostic, not as the definition of
collaborative quality.

\paragraph{D4: head-tail capacity.}
\tool splits items into popularity buckets and measures whether full \sid
capacity and prefix structure are allocated differently across head, mid, and
tail items. This helps separate genuine tokenizer behavior from simply
restating the catalog popularity distribution.

\paragraph{D5a and D6.}
D5a reports structural cost proxies: \sid length, prefix fan-out, duplicate
codes, and trie-like expansion pressure. These are not real generator serving
latencies. If generator outputs are available, future D5b/D7 metrics can add
invalid paths, next-token entropy, and generated-candidate duplication. D6
computes churn between tokenizer refreshes and is useful for drift-aware
continual-tokenization settings~\cite{feng2026dact}.

\begin{table*}[t]
\caption{Method coverage and artifact role in the current \tool resource.}
\label{tab:method-coverage}
\small
\begin{tabularx}{\textwidth}{@{}llllY@{}}
\toprule
Line & Cluster & Dataset & Artifact & Paper role and caveat \\
\midrule
GRID/RQ-KMeans & A & All\_Beauty; Musical & real export &
Main A evidence; Musical row uses processed feature text, not raw-text TIGER. \\
ReSID/GAOQ & B & Musical & real export &
Main B evidence; bounded one-epoch FAMAE and GAOQ export. \\
Sanity tokenizers & control & Musical; MovieLens smoke & deterministic &
Metric non-redundancy controls; not named methods. \\
CARD compact & B/control & Musical; Sports & proxy/stressor &
Backlog or stressor only; not faithful CARD evidence. \\
DACT & drift & Tools & bundled arrays &
Optional D6 churn demo; not a Cluster B replacement. \\
DIGER/CapsID/AdaSID/AsymRec & future & not run & not claimed &
Coverage context only; verify releases before submission. \\
\bottomrule
\end{tabularx}
\end{table*}
